\begin{table}[htbp]
\centering

\small
\begin{tabular}{lccc}
\toprule
 & (1) & (2) & (3) \\
 & Composite & Predictions vs. & Re-estimated on \\
Dimension & (headline) & visit-only & visit-only \\
\midrule
Warmth     & $0.765$ & $0.267$ & $0.228$ \\
           & ($p<0.001$) & ($p=0.007$) & ($p=0.022$) \\
\addlinespace
Strictness & $0.701$ & $0.276$ & $0.349$ \\
           & ($p<0.001$) & ($p=0.005$) & ($p<0.001$) \\
\bottomrule
\multicolumn{4}{l}{\footnotesize Entries are Pearson $r$; leave-one-out
cross-validated throughout.}\\
\end{tabular}
\caption{Interview-tier model performance, decomposed}
\label{tab:ridge_circularity}
\begin{minipage}{\linewidth}
\smallskip
\footnotesize\textit{Notes:} Column~(1) is the headline figure: leave-one-out predictions correlated with the gold-standard composite (the combined visit-and-interview culture score), of which the interview sub-scores form 40~per~cent by construction. Column~(2) correlates those same predictions with the visit-only component, which shares no input with the predictors. Column~(3) re-estimates the identical feature set directly against the visit-only component, with the ridge penalty re-selected by generalised cross-validation for that target, giving the best that these features can do against an independent target. $N = 101$ full-data schools with complete features. Reproduced by \texttt{ridge\_diagnostics.py}.
\end{minipage}
\end{table}
